\section{实验结果}
\label{sec:results}

\subsection{整数正确性与系统稳定性}

\cref{tab:correctness} 汇总当前已完成正确性证据。128 张代表图共比较
\ConformanceImages{}$\times$22=\NodeRecords{} 个节点记录，累计约 1.083 GB
张量，board 与 RTL-semantic host 的 payload SHA 完全相同，逐字节 mismatch
为 \NodeMismatches{}。双 raw head 的 128 条记录同样完全一致。完整 5000 张
val2017 raw-head 经 canonical host decode/NMS 后，12 项 COCOeval 统计和
80 类 AP 与 host INT8 的差值均为 0。这证明当前 bitstream 执行的网络整数
语义与 host 权威模型一致。板端 A53 accuracy product 也完成全部
\ProductImages{} 张：
逐图检测数、类别、source index、输出顺序和 NMS 保留集合完全一致，score
最大绝对差为 $\ProductScoreDelta$，bbox 最大绝对差为
$\ProductBBoxDelta$ 像素，12 项 COCO 指标最大差为 \ProductMetricDelta{}。

\begin{table*}[t]
\centering
\caption{当前可引用的正确性和稳定性证据。}
\label{tab:correctness}
\small
\resizebox{0.98\textwidth}{!}{%
\begin{tabular}{lrrrrl}
\toprule
项目 & 图像/记录 & 节点/头 & 比较字节 & 最大差异 & 状态\\
\midrule
逐节点 RTL-semantic conformance & 128 & 2816节点 & 1,082,530,176 & 0 byte & PASS \\
128张双 raw head & 128 & 256头 & 27,580,800 & 0 byte & PASS \\
5000张 network mAP & 5000 & 10000头 & -- & 12项/每类差0 & PASS \\
5000张 A53 accuracy product & 5000 & detections & -- & score $5.96\!\times\!10^{-8}$；bbox $1.37\!\times\!10^{-4}$ px & PASS \\
600 s Ethernet soak & 13,184记录 & -- & -- & CRC/协议/DMA/PL/timeout=0 & PASS \\
\bottomrule
\end{tabular}
}
\end{table*}


600 s soak 实际持续 \SoakSeconds{} s，循环 128 张集合产生 13,184 条记录，
输出 CRC mismatch、协议错误和意外 session reconnect 均为 0，记录到的最高
PS 温度为 \SoakTempC{}\,$^\circ$C。该测试说明确定输入下输出稳定，并覆盖
持久网络会话、重复 DMA 和长时间 cache 操作；它不替代电源循环或环境温箱
可靠性测试。

\begin{figure*}[t]
\centering
\begin{tikzpicture}[font=\small,
  done/.style={draw,rounded corners,fill=LasaGreen!18,minimum width=31mm,minimum height=10mm,align=center},
  wait/.style={draw,rounded corners,fill=LasaPending,minimum width=31mm,minimum height=10mm,align=center},
  arr/.style={->,thick,LasaGray}]
\node[done] (rtl) {RTL/形式化与回归\\已完成};
\node[done,right=7mm of rtl] (impl) {200 MHz实现门禁\\已完成};
\node[done,right=7mm of impl] (exact) {128张逐节点\\byte-exact};
\node[done,right=7mm of exact] (raw) {5000张raw-head\\mAP一致};
\node[done,below=9mm of rtl] (perf) {3$\times$1000性能\\已完成};
\node[done,right=7mm of perf] (soak) {600 s稳定性\\已完成};
\node[done,right=7mm of soak] (prod) {5000张product\\PASS};
\node[wait,right=7mm of prod] (abl) {消融/CPU/自检\\待补};
\draw[arr] (rtl)--(impl); \draw[arr] (impl)--(exact); \draw[arr] (exact)--(raw);
\draw[arr] (perf)--(soak); \draw[arr] (soak)--(prod); \draw[arr] (prod)--(abl);
\node[draw,dashed,fit=(rtl)(impl)(exact)(raw)(perf)(soak)(prod),inner sep=3mm,label=below:{当前可证实论文事实}] {};
\node[draw,dashed,LasaOrange,fit=(abl),inner sep=3mm,label=below:{投稿冻结硬门禁}] {};
\end{tikzpicture}
\caption{论文证据链与冻结状态。绿色项目已由文件和哈希闭环；黄色项目
仍需真实实验，当前初稿只给出方法和表格位置，不提前填入结果。}
\label{fig:evidence-matrix}
\end{figure*}


\subsection{模型精度}

\begin{table*}[t]
\centering
\caption{COCO val2017 精度（百分数）。640 结果仅用于上游资产复现；416 结果才用于部署公平比较。}
\label{tab:accuracy}
\small
\begin{tabular}{lrrrrrr}
\toprule
配置 & AP & AP50 & AP75 & AP$_S$ & AP$_M$ & AP$_L$ \\
\midrule
官方 FP32-640 & 17.636 & 34.862 & -- & -- & -- & -- \\
部署 FP32-416 & 17.494 & 33.493 & 16.415 & 5.324 & 18.436 & 29.263 \\
当前部署 INT8-416 & 14.304 & 30.212 & 12.145 & 4.325 & 15.245 & 23.686 \\
\bottomrule
\end{tabular}
\end{table*}


官方 640 基线 AP50:95/AP50 为 17.636\%/34.862\%，与上游公开量级一致；
固定 416 预处理下 FP32 为 17.494\%/33.493\%。当前部署 epoch-1 INT8 为
14.304\%/30.212\%，相对部署 FP32 分别下降 3.190 和 3.281 个百分点，超过
预设的 1.0/2.0 发布门限。AP75、small、medium 和 large 同样下降，表明损失
不是只由某个阈值或目标尺度造成。

由于 5000 张板端 raw head 与 RTL host 完全一致，board network mAP 必然等于
当前 host INT8 数值；因此继续修改 DMA、调度或 decode 不能提高 network mAP。
模型精度需要更完整的服务器 QAT、校准和 checkpoint 选择。本文把 epoch-1
作为功能合格参数，用于证明硬件与系统语义；不会使用“无精度损失”或“模型
发布通过”的表述。未来替换参数包时需要重新执行 128 张 byte-exact、5000 张
network/product 和性能/soak，而不需要改变 \LASA 架构结论。

\subsection{端到端性能}

\begin{table}[t]
\centering
\caption{板端主性能结果，3 次独立运行、每次 20 次预热和 1000 次计时。}
\label{tab:performance}
\small
\begin{tabular}{lrrl}
\toprule
指标 & 均值/值 & P95 & 单位 \\
\midrule
Resident 总延迟 & 34.943 & 34.965 & ms \\
PL 13 卷积 & 33.607 & -- & ms \\
A53（含 decode） & 1.301 & -- & ms \\
双头 decode/NMS & 0.175 & -- & ms \\
Resident 吞吐 & 28.618 & -- & FPS \\
网络 pipeline 吞吐 & 21.263 & -- & FPS \\
\bottomrule
\end{tabular}
\end{table}


正式 3$\times$1000 结果显示 resident 平均为 \ResidentMeanMs{} ms，P95 为
\ResidentPNinetyFiveMs{} ms，对应 \ResidentFPS{} FPS；包含当前 TCP 分块开销的
pipeline 为 21.263 FPS。PL 13 个卷积平均 \PLMeanMs{} ms，占 resident 的
96.2\%；A53 含 decode 平均 \AFTMeanMs{} ms，其中 decode/NMS 约
\DecodeMeanMs{} ms。三次运行的 3000 条输出 CRC 序列一致，说明性能统计不是
来自错误或跳过的图像。

\begin{figure*}[t]
\centering
\begin{tikzpicture}
\begin{axis}[
  xbar stacked,width=0.90\textwidth,height=4.8cm,
  xmin=0,xmax=36,xlabel={平均延迟/ms},
  ymin=-0.6,ymax=0.6,ytick=data,yticklabels={完整 resident},
  legend style={at={(0.5,1.04)},anchor=south,legend columns=4,draw=none,font=\small},
]
\addplot[fill=LasaBlue!75] coordinates {(\PLMeanMs,0)}; \addlegendentry{PL 13卷积}
\addplot[fill=LasaOrange!75] coordinates {(\AFTensorMeanMs,0)}; \addlegendentry{A53张量算子/调度}
\addplot[fill=LasaGreen!70] coordinates {(\DecodeMeanMs,0)}; \addlegendentry{decode/NMS}
\addplot[fill=gray!45] coordinates {(\ResidualMeanMs,0)}; \addlegendentry{残余计时差}
\end{axis}
\end{tikzpicture}
\par\smallskip
\parbox{0.90\textwidth}{\centering\small
m13：\MThirteenMeanMs{} ms；m19：\MNineteenMeanMs{} ms；
两者合计占 PL 时间 \LongLayerPLShare{}\%。}
\caption{正式 3$\times$1000 resident 延迟分解。A53 含 decode 的均值为
\AFTMeanMs{} ms，说明当前端到端延迟主要由 PL 长卷积层决定。}
\label{fig:breakdown}
\end{figure*}


模型总量为 2.782480896 GMAC（5.564961792 GOP）。按 PL 平均时间，系统完成
13 个卷积的有效吞吐约 \PLEffectiveGOPS{} GOPS；相对 576 MAC/cycle、
200 MHz 的 230.4 GOPS 理论峰值，利用率约 \ArrayUtilization{}\%。该利用率
包含尾 lane、Cout 尾块、权重/输入、PSUM feedback 和上下文空泡，是比只在
阵列 fire 周期计算更严格的网络有效利用率。

\begin{table}[t]
\centering
\caption{13 个 PL 卷积的平均延迟及其占 PL 总时间的比例。}
\label{tab:layer-latency}
\small
\begin{tabular}{lrr}
\toprule
层 & 延迟/ms & 占PL/\% \\
\midrule
m0 & 3.015 & 8.97 \\
m2 & 2.253 & 6.70 \\
m4 & 2.120 & 6.31 \\
m6 & 2.065 & 6.14 \\
m8 & 2.135 & 6.35 \\
m10 & 2.197 & 6.54 \\
m13 & 8.589 & 25.56 \\
m14 & 1.020 & 3.04 \\
m15 & 2.197 & 6.54 \\
m16 & 0.307 & 0.91 \\
m19 & 6.240 & 18.57 \\
p4\_detect & 0.875 & 2.60 \\
p5\_detect & 0.593 & 1.77 \\
\bottomrule
\end{tabular}
\end{table}


m13 与 m19 平均分别为 8.589 ms 和 6.240 ms，合计约占 PL 时间 44.1\%。
m13 受 256 个 $K$-pass 和 32 个 Cout block 影响，m19 同时具有 192 pass、
26$\times$26 空间和 5 个 spatial tile。检测头虽为 255 通道，但原生
1$\times$1 使 P4/P5 仅约 0.875/0.593 ms，说明尾通道支持没有把检测头变成
主要瓶颈。下一步硬件优化应优先减少长 3$\times$3 层的 pass/context 开销，
而不是继续微调已经只有 0.175 ms 的 decode。

SD 冷启动后的 20 warmup+100 timed 复测得到 resident 34.925 ms、PL 33.602 ms、
A53 1.294 ms 和 28.632 FPS，与正式 JTAG/网络驻留结果接近。该结果支持冷启动
路径与 JTAG 下载使用同一计算内核，但样本数较小，只作为重复性证据，不能
替代 3000 条主结果。

\subsection{受控消融与 CPU 基线}

\cref{tab:ablation-plan} 给出冻结表结构。当前版本尚未生成所有受控变体的
正式 200 MHz bitstream，因此本节不使用历史 100 MHz、18$\times$8 或单尺度
结果填空。尤其不能把不同 RTL 根、不同数组规模或不能执行完整网络的早期
结果直接相减，称为某一个机制的加速。

\begin{table*}[t]
\centering
\caption{受控消融结果表（当前为冻结前执行合同；数值必须由同源 200 MHz 变体自动回填）。}
\label{tab:ablation-plan}
\small
\resizebox{0.98\textwidth}{!}{%
\begin{tabular}{lrrrrrrl}
\toprule
变体 & 延迟/ms & 相对加速 & DDR字节 & compute idle & context gap & 资源增量 & byte-exact\\
\midrule
稀疏3$\times$3模拟1$\times$1 & PENDING & PENDING & PENDING & PENDING & PENDING & PENDING & 必须PASS \\
原生1$\times$1 & PENDING & PENDING & PENDING & PENDING & PENDING & PENDING & 必须PASS \\
逐pass IFM & PENDING & PENDING & PENDING & PENDING & PENDING & PENDING & 必须PASS \\
HWC物化/重放 & PENDING & PENDING & PENDING & PENDING & PENDING & PENDING & 必须PASS \\
无预载/无handoff & PENDING & PENDING & PENDING & PENDING & PENDING & PENDING & 必须PASS \\
仅预载 & PENDING & PENDING & PENDING & PENDING & PENDING & PENDING & 必须PASS \\
预载+handoff & PENDING & PENDING & PENDING & PENDING & PENDING & PENDING & 必须PASS \\
完整跨pass重叠 & PENDING & PENDING & PENDING & PENDING & PENDING & PENDING & 必须PASS \\
固定保守tile & PENDING & PENDING & PENDING & PENDING & PENDING & PENDING & 必须PASS \\
层自适应tile（本文） & 34.943 & 1.000 & PENDING & PENDING & PENDING & -- & PASS \\
\bottomrule
\end{tabular}
}
\end{table*}


同板 CPU 单核标量与四核 NEON 基线同样处于 PENDING。它们的意义不是证明
FPGA 必然优于任意 CPU 库，而是隔离相同整数语义和相同数据布局下的计算
收益。投稿版必须报告完整 13 卷积和相同 A53 后处理，不能只挑 m13 或 m19
推算网络加速。

\subsection{资源、时序与功耗估算}

\begin{table*}[t]
\centering
\caption{r5 最终系统 post-route 资源、时序与功耗估算。}
\label{tab:resources}
\small
\begin{tabular}{lrr}
\toprule
指标 & 使用/结果 & 备注\\
\midrule
LUT & 56,949 & logic 51,152；LUTRAM 5,797 \\
CLB & 11,146/14,640 & 76.13\% \\
BRAM & 94 & 65.28\% \\
URAM & 48 & 75.00\% \\
DSP & 650 & 阵列 576 \\
WNS & 0.004 ns & setup EP=0 \\
WHS & 0.009 ns & hold EP=0 \\
Route/DRC Error & 0/0 & fully routed \\
Post-route estimated power & 4.008 W & dynamic 3.695 W；static 0.313 W \\
Vectorless PL/PS dynamic & 1.133/2.562 W & accelerator hierarchy 0.996 W \\
Estimated junction/confidence & 34.3 $^{\circ}$C / Medium & ambient 25.0 $^{\circ}$C \\
\bottomrule
\end{tabular}
\end{table*}


最终系统使用 \LUTCount{} LUT、\BRAMCount{} BRAM、\URAMCount{} URAM 和
\DSPCount{} DSP；CLB site 占用 76.13\%。200 MHz route WNS 只有
\WNSValue{} ns，WHS \WHSValue{} ns，所有 setup/hold/pulse failing endpoint
为零。4 ps setup 余量说明设计满足当前 formal gate，但几乎没有进一步加深
关键控制锥的空间。URAM 已用 75\%，扩展物化容量也受到明显约束。因此后续
微架构应采用局部流水和物理友好复制，避免在 context/descriptor 共享网络上
增加组合级数。

正式 routed DCP 的 vectorless post-route 功耗估算为 \PostRoutePowerTotalW{} W，
其中动态功耗 \PostRoutePowerDynamicW{} W、器件静态功耗
\PostRoutePowerStaticW{} W。PS8 动态功耗为 \PostRoutePowerPSDynamicW{} W；
按总动态功耗扣除 PS8 得到 PL 动态功耗 \PostRoutePowerPLDynamicW{} W，
其中加速器层次为 \PostRoutePowerAccelDynamicW{} W。估算采用 typical process、
25\,$^\circ$C 环境温度、12.5\% vectorless toggle rate 和 0.5 static probability，
未加载 SAIF；估算结温为 \PostRoutePowerJunctionC{}\,$^\circ$C，活动率置信度为
\PostRoutePowerConfidence{}。因此这些数值用于识别功耗构成，不能视为板级实测，
也不与外部工作的实测 GOPS/W 直接比较。

\subsection{外部工作对比}

\begin{table*}[t]
\centering
\caption{外部工作的比较口径。当前初稿只冻结能力与字段；投稿前由原始论文自动生成数值，避免跨口径误排。}
\label{tab:external-comparison}
\small
\resizebox{0.98\textwidth}{!}{%
\begin{tabular}{lcccccccl}
\toprule
工作 & 器件 & MHz & 网络/输入 & 完整双头 & 后处理 & batch & 延迟/FPS & 口径备注\\
\midrule
YOLOv2 FPGA\cite{yolov2fpga2018} & 见原文 & 见原文 & YOLOv2 & -- & 见原文 & 见原文 & 见原文 & 非同网络 \\
YOLOv3 FPGA\cite{yolov3fpga2019} & 见原文 & 见原文 & YOLOv3 & 见原文 & 见原文 & 见原文 & 见原文 & 核对完整性 \\
Tiny-YOLOv3 accelerator\cite{tinyyoloaccel2022} & 见原文 & 见原文 & tiny/见原文 & 见原文 & 见原文 & 见原文 & 见原文 & 原文口径 \\
Multicore YOLOv3-tiny\cite{multicoreyolo2023} & 见原文 & 见原文 & tiny/见原文 & 见原文 & 见原文 & 见原文 & 见原文 & 直接对照 \\
本文 \LASA & XCK26 & 200 & tiny/416 & 是 & A53包含 & 1 & 34.943 ms/28.618 FPS & resident \\
\bottomrule
\end{tabular}
}
\end{table*}


当前表中的“见原文”是显式冻结占位，不是缺失值按零处理。投稿版需记录每个
外部数值的页码/表号、是否 conv-only、是否包含 CPU 后处理、输入和 batch；
若无法统一口径，保留原文数字但不据此宣称绝对领先。

\begin{table}[t]
\centering
\caption{投稿冻结门禁。PENDING 项不得在摘要、结论或对比表中写成已完成结果。}
\label{tab:freeze-gates}
\footnotesize
\begin{tabular}{lll}
\toprule
实验/材料 & 状态 & 处理 \\
\midrule
5000 张 A53 product 后处理 & PASS & 投稿冻结前硬门禁 \\
受控微架构消融 & PENDING & 投稿冻结前硬门禁 \\
A53 单核标量基线 & PENDING & 投稿冻结前硬门禁 \\
A53 四核 NEON 基线 & PENDING & 投稿冻结前硬门禁 \\
Post-route 功耗估算 & PASS & 投稿冻结前硬门禁 \\
公开 artifact 自检 & PENDING & 投稿冻结前硬门禁 \\
\bottomrule
\end{tabular}
\end{table}

